% id: 50-08
% book: 베이직쎈 미적분1 · 03 미분계수와 도함수
% section: 유형 022 평균변화율과 미분계수의 기하적 의미
% figure: none
% answer: 
% answer_source: 
% variant_level: 0 (원본 전사)
% 이 파일은 build.mjs 가 items.json 에서 만든다. 고칠 때는 items.json 을 고친다.
\smash{\raisebox{1.5mm}{\dmkichul{베이직쎈 미적분1 50쪽 08번}}}
함수 $y=f(x)$의 그래프는 두 점~$\pt{P}(1{,}\mkern3mu\mathopen{}f(1))$, $\pt{Q}(4{,}\mkern3mu\mathopen{}f(4))$를 지난다. 두 점~$\pt{P}$, $\pt{Q}$를 지나는 직선의 기울기가 $-2$일 때, 함수 $f(x)$에서 $x$의 값이 $1$에서 $4$까지 변할 때의 평균변화율을 구하시오.
